% GENERATED FILE. Edit canonical lessons, then run npm run generate:books.
% canonical-source-hash: fnv1a64:ec3d318a62743e09
% canonical-lessons: TE-C330-vilekari, TE-C330-patrika, TE-C330-sirsika, TE-C330-kathanam, TE-C330-pracurincu

\chapter{The Reporter}
\label{ch:te-the-reporter}
\hlchaptermodality{330}

\begin{chapteropening}
\textbf{By the end of this chapter:} \emph{I can say a reporter, a periodical, a headline, a news story and to publish.}

Complete the last lesson of chapter 330: I can say a reporter, a periodical, a headline, a news story and to publish.
\end{chapteropening}

\section[\texorpdfstring{\te{విలేకరి} (vilēkari)}{విలేకరి (vilēkari)}]{\te{విలేకరి} (vilēkari) — a reporter; a correspondent}

\label{lesson:TE-C330-vilekari}



\begin{quote}
Before the new one: say the Telugu for to make off with, then the Telugu for to sense.
\end{quote}

\subsection*{You'll want to know: \te{విలేకరి}}
\textbf{\te{విలేకరి}} — \emph{vilēkari} — \textquotedblleft{}a reporter; a correspondent\textquotedblright{}.

\begin{quote}
\te{విలేకరి} \te{సాక్షితో} \te{మాట్లాడాడు}.
\end{quote}

\begin{quote}
\emph{vilēkari sākṣitō māṭlāḍāḍu.}
\end{quote}

\begin{quote}
The reporter spoke to the witness.
\end{quote}

Plural: \te{విలేకరులు}, \emph{vilēkarulu}.

A paper's reporter in a town is its \te{విలేకరి} there.

\begin{cousinweb}
\textbf{\te{విలేకరి}} (\emph{vilēkari}) is the Telugu word for a newspaper reporter; its origin is not settled.
\end{cousinweb}

\begin{grammarlens}[title={What you've built}]
One more word for telling what happened.
\end{grammarlens}

\subsection*{Your turn}
\begin{itemize}
\raggedright
  \item \emph{Say it:} \emph{vilēkari}
  \item \emph{Say it:} \emph{vilēkari}, once more
  \item \emph{Say it:} \emph{vilēkari sākṣitō māṭlāḍāḍu.} — The reporter spoke to the witness.
  \item \emph{Recall:} say the Telugu for to make off with, then the Telugu for to sense, then say \emph{vilēkari} again
\end{itemize}

\begin{tcolorbox}[breakable,colback=teal!4,colframe=teal!35!black,title={Before you move on}]
What does \te{విలేకరి} mean? (\textquotedblleft{}A reporter; a correspondent\textquotedblright{}.) Say it once more.
\end{tcolorbox}
\section[\texorpdfstring{\te{పత్రిక} (patrika)}{పత్రిక (patrika)}]{\te{పత్రిక} (patrika) — a periodical; a magazine}

\label{lesson:TE-C330-patrika}



\begin{quote}
Before the new one: say the Telugu for to sense, then the Telugu for a reporter.
\end{quote}

\subsection*{You'll want to know: \te{పత్రిక}}
\textbf{\te{పత్రిక}} — \emph{patrika} — \textquotedblleft{}a periodical; a magazine\textquotedblright{}.

\begin{quote}
\te{ఆ} \te{పత్రిక} \te{ప్రతి} \te{వారం} \te{వస్తుంది}.
\end{quote}

\begin{quote}
\emph{ā patrika prati vāraṁ vastundi.}
\end{quote}

\begin{quote}
That magazine comes out every week.
\end{quote}

Plural: \te{పత్రికలు}, \emph{patrikalu}.

It is the second half of \te{వార్తాపత్రిక}, a newspaper: a news sheet.

\begin{cousinweb}
\textbf{\te{పత్రిక}} (\emph{patrika}) is borrowed from Sanskrit \emph{pattrikā}, \textquotedblleft{}a leaf for writing upon; a letter, a document\textquotedblright{} (Monier-Williams).
\end{cousinweb}

\begin{grammarlens}[title={What you've built}]
One more word for telling what happened.
\end{grammarlens}

\subsection*{Your turn}
\begin{itemize}
\raggedright
  \item \emph{Say it:} \emph{patrika}
  \item \emph{Say it:} \emph{patrika}, once more
  \item \emph{Say it:} \emph{ā patrika prati vāraṁ vastundi.} — That magazine comes out every week.
  \item \emph{Recall:} say the Telugu for to sense, then the Telugu for a reporter, then say \emph{patrika} again
\end{itemize}

\begin{tcolorbox}[breakable,colback=teal!4,colframe=teal!35!black,title={Before you move on}]
What does \te{పత్రిక} mean? (\textquotedblleft{}A periodical; a magazine\textquotedblright{}.) Say it once more.
\end{tcolorbox}
\section[\texorpdfstring{\te{శీర్షిక} (śīrṣika)}{శీర్షిక (śīrṣika)}]{\te{శీర్షిక} (śīrṣika) — a headline; a title}

\label{lesson:TE-C330-sirsika}



\begin{quote}
Before the new one: say the Telugu for a reporter, then the Telugu for a periodical.
\end{quote}

\subsection*{You'll want to know: \te{శీర్షిక}}
\textbf{\te{శీర్షిక}} — \emph{śīrṣika} — \textquotedblleft{}a headline; a title\textquotedblright{}.

\begin{quote}
\te{పత్రికలో} \te{పెద్ద} \te{శీర్షిక} \te{వచ్చింది}.
\end{quote}

\begin{quote}
\emph{patrikalō pedda śīrṣika vaccindi.}
\end{quote}

\begin{quote}
A big headline appeared in the paper.
\end{quote}

The line at the head of a page or a story.

\begin{cousinweb}
\textbf{\te{శీర్షిక}} (\emph{śīrṣika}) is a learned word built on Sanskrit \emph{śīrṣa}, \textquotedblleft{}the head\textquotedblright{} (CDIAL 12497): the words at the head of a story.
\end{cousinweb}

\begin{grammarlens}[title={What you've built}]
One more word for telling what happened.
\end{grammarlens}

\subsection*{Your turn}
\begin{itemize}
\raggedright
  \item \emph{Say it:} \emph{śīrṣika}
  \item \emph{Say it:} \emph{śīrṣika}, once more
  \item \emph{Say it:} \emph{patrikalō pedda śīrṣika vaccindi.} — A big headline appeared in the paper.
  \item \emph{Recall:} say the Telugu for a reporter, then the Telugu for a periodical, then say \emph{śīrṣika} again
\end{itemize}

\begin{tcolorbox}[breakable,colback=teal!4,colframe=teal!35!black,title={Before you move on}]
What does \te{శీర్షిక} mean? (\textquotedblleft{}A headline; a title\textquotedblright{}.) Say it once more.
\end{tcolorbox}
\section[\texorpdfstring{\te{కథనం} (kathanaṁ)}{కథనం (kathanaṁ)}]{\te{కథనం} (kathanaṁ) — a news story; an account}

\label{lesson:TE-C330-kathanam}



\begin{quote}
Before the new one: say the Telugu for a periodical, then the Telugu for a headline.
\end{quote}

\subsection*{You'll want to know: \te{కథనం}}
\textbf{\te{కథనం}} — \emph{kathanaṁ} — \textquotedblleft{}a news story; an account\textquotedblright{}.

\begin{quote}
\te{విలేకరి} \te{ఆ} \te{సంఘటన} \te{మీద} \te{కథనం} \te{రాశాడు}.
\end{quote}

\begin{quote}
\emph{vilēkari ā saṅghaṭana mīda kathanaṁ rāśāḍu.}
\end{quote}

\begin{quote}
The reporter wrote a story on that incident.
\end{quote}

A cousin of \te{కథ}, a story: the news calls each report a \te{కథనం}.

\begin{cousinweb}
\textbf{\te{కథనం}} (\emph{kathanaṁ}) is borrowed from Sanskrit \emph{kathana}, \textquotedblleft{}narration, telling\textquotedblright{} (CDIAL 2702), from the root \emph{kath}, \textquotedblleft{}to tell\textquotedblright{}, which also gives \te{కథ}, a story.
\end{cousinweb}

\begin{grammarlens}[title={What you've built}]
One more word for telling what happened.
\end{grammarlens}

\subsection*{Your turn}
\begin{itemize}
\raggedright
  \item \emph{Say it:} \emph{kathanaṁ}
  \item \emph{Say it:} \emph{kathanaṁ}, once more
  \item \emph{Say it:} \emph{vilēkari ā saṅghaṭana mīda kathanaṁ rāśāḍu.} — The reporter wrote a story on that incident.
  \item \emph{Recall:} say the Telugu for a periodical, then the Telugu for a headline, then say \emph{kathanaṁ} again
\end{itemize}

\begin{tcolorbox}[breakable,colback=teal!4,colframe=teal!35!black,title={Before you move on}]
What does \te{కథనం} mean? (\textquotedblleft{}A news story; an account\textquotedblright{}.) Say it once more.
\end{tcolorbox}
\section[\texorpdfstring{\te{ప్రచురించు} (pracuriñcu)}{ప్రచురించు (pracuriñcu)}]{\te{ప్రచురించు} (pracuriñcu) — to publish}

\label{lesson:TE-C330-pracurincu}



\begin{quote}
Before the new one: say the Telugu for a headline, then the Telugu for a news story.
\end{quote}

\subsection*{You'll want to know: \te{ప్రచురించు}}
\textbf{\te{ప్రచురించు}} — \emph{pracuriñcu} — \textquotedblleft{}to publish\textquotedblright{}.

\begin{quote}
\te{పత్రిక} \te{ఆ} \te{కథనం} \te{ప్రచురించింది}.
\end{quote}

\begin{quote}
\emph{patrika ā kathanaṁ pracuriñcindi.}
\end{quote}

\begin{quote}
The magazine published that story.
\end{quote}

Past (I …): \te{ప్రచురించాను}, \emph{pracuriñcānu}.

A -\te{ఇంచు} verb: Telugu makes verbs of Sanskrit words with -\te{ఇంచు}, as in \te{ముద్రించు}, to print.

\begin{cousinweb}
\textbf{\te{ప్రచురించు}} (\emph{pracuriñcu}) is built on Sanskrit \emph{pracura}, \textquotedblleft{}abundant, plentiful\textquotedblright{} (Monier-Williams), with the Telugu verb ending \emph{-iñcu}: to publish is to spread a text wide.
\end{cousinweb}

\begin{grammarlens}[title={What you've built}]
One more word for telling what happened.
\end{grammarlens}

\subsection*{Your turn}
\begin{itemize}
\raggedright
  \item \emph{Say it:} \emph{pracuriñcu}
  \item \emph{Say it:} \emph{pracuriñcu}, once more
  \item \emph{Say it:} \emph{patrika ā kathanaṁ pracuriñcindi.} — The magazine published that story.
  \item \emph{Recall:} say the Telugu for a headline, then the Telugu for a news story, then say \emph{pracuriñcu} again
\end{itemize}

\begin{tcolorbox}[breakable,colback=teal!4,colframe=teal!35!black,title={Before you move on}]
What does \te{ప్రచురించు} mean? (\textquotedblleft{}To publish\textquotedblright{}.) Say it once more.
\end{tcolorbox}
